% ====================================================================
% ch1_preamble.tex — Chapter 1 (v1.0.21) 공통 preamble·매크로
%   graphite_ica_ch1_v1.0.21.tex 가 \input 한다. 절별 파일은 _sections/ 에.
%   빌드: xelatex -interaction=nonstopmode (kotex — 한글 전제), 3-pass.
%   렌더 목표 = v1.0.18.2 동등(문서클래스·패키지·박스 체계·매크로 계승),
%   조립만 절별 \input 으로 전환.
% ====================================================================
\documentclass[11pt,a4paper]{article}
\usepackage[margin=22mm]{geometry}
\usepackage{kotex}
\setmonofont{D2Coding}[Scale=MatchLowercase]
\usepackage{amsmath,amssymb,mathtools,bm}
\usepackage{booktabs,longtable,array}
\usepackage{enumitem}
\usepackage{amsthm}
\usepackage{hyperref}
\usepackage{fancyhdr}
\usepackage{lastpage}
\usepackage{setspace}
\usepackage{tikz}
\usetikzlibrary{positioning,arrows.meta,calc,fit,backgrounds}
\setstretch{1.12}
\setlength{\parskip}{0.45em}\setlength{\parindent}{0pt}\setlength{\headheight}{14pt}
\setlength{\emergencystretch}{2em}
\hypersetup{colorlinks=true, linkcolor=blue, urlcolor=blue, citecolor=blue,
  pdftitle={흑연 음극 + LCO 양극 dQ/dV 이론: 계산 진행을 따라가는 수식-구동 전개 — Chapter 1 (v1.0.21)}}
\pagestyle{fancy}\fancyhf{}
\lhead{흑연 음극 + LCO 양극 $dQ/dV$ 이론 — Ch.1 (v1.0.21, 통계역학-first + 계산 진행)}\rhead{\thepage/\pageref{LastPage}}
\renewcommand{\headrulewidth}{0.4pt}

\newtheorem*{keybox}{핵심}
\newtheorem*{codebox}{코드 대응}
\newtheorem*{signbox}{부호 검산}
\newtheorem*{verifybox}{검산}
\newtheorem*{derivbox}{유도}
\newtheorem*{bgbox}{배경}

\newcommand{\dd}{\mathrm{d}}
\newcommand{\eff}{\mathrm{eff}}
\newcommand{\eq}{\mathrm{eq}}
\newcommand{\app}{\mathrm{app}}
\newcommand{\cell}{\mathrm{cell}}
\newcommand{\bg}{\mathrm{bg}}
\newcommand{\rxn}{\mathrm{rxn}}
\newcommand{\dis}{\mathrm{dis}}
\newcommand{\chg}{\mathrm{chg}}
\newcommand{\hys}{\mathrm{hys}}
\newcommand{\work}{\mathrm{work}}
\newcommand{\rep}{\mathrm{rep}}
\newcommand{\cat}{\mathrm{cat}}
\newcommand{\code}[1]{\texttt{#1}}

\renewcommand{\theequation}{1.\arabic{equation}}
\renewcommand{\thesection}{1.\arabic{section}}

\makeatletter
\renewcommand*\l@section[2]{%
  \ifnum \c@tocdepth >\z@
    \addpenalty\@secpenalty
    \addvspace{1.0em \@plus\p@}%
    \setlength\@tempdima{2.3em}%
    \begingroup
      \parindent \z@ \rightskip \@pnumwidth
      \parfillskip -\@pnumwidth
      \leavevmode \bfseries
      \advance\leftskip\@tempdima
      \hskip -\leftskip
      #1\nobreak\hfil
      \nobreak\hb@xt@\@pnumwidth{\hss #2\kern-\p@\kern\p@}\par
    \endgroup
  \fi}
\renewcommand*\l@subsection{\@dottedtocline{2}{2.3em}{3.2em}}
\makeatother
